\documentclass{article}
\usepackage{graphicx}
\usepackage{hyperref}
\usepackage{booktabs}
\usepackage{geometry}

\geometry{a4paper, margin=1in}

\title{Predicting Breast Cancer Progression Risk Using Biologically-Informed Feature Engineering and Explainable Machine Learning}
\author{Suhaan Thayyil \and Eshaan Nidee}
\date{ISEF 2026}

\begin{document}

\maketitle

\begin{abstract}
Breast cancer prognosis varies significantly even among patients with similar clinical features. This study leverages machine learning on gene expression data to predict progression risk. By engineering biologically interpretable features (pathway activity scores) and utilizing Explainable AI (SHAP), we developed a model that achieves an AUC of 0.86 on an independent external validation cohort (GSE96058). Our findings confirm that proliferation and estrogen response pathways are critical drivers of risk, offering a transparent and accurate tool for clinical decision support.
\end{abstract}

\section{Introduction}
Breast cancer is a heterogeneous disease. While molecular subtypes have improved stratification, substantial intra-subtype variation in prognosis remains. Traditional clinical factors often fail to capture the underlying biological drivers of aggression. High-throughput gene expression profiling offers a solution, but "black box" machine learning models often lack the interpretability required for clinical adoption. This project aims to bridge that gap by using biologically informed feature engineering and SHAP analysis.

\section{Methods}

\subsection{Data Sources}
We utilized two large-scale public datasets:
\begin{itemize}
    \item \textbf{Training (TCGA-BRCA):} $N=213$ primary tumor samples with RNA-seq expression and complete survival outcome data.
    \item \textbf{Validation (GSE96058):} $N=1,483$ samples from the SCAN-B cohort, serving as an independent external validation set.
\end{itemize}

\subsection{Preprocessing \& Feature Engineering}
Gene expression data were log2-transformed and z-score normalized. Instead of using raw individual gene expression values, which can be noisy, we aggregated genes into biologically meaningful \textbf{Pathway Scores} using Single-Sample Gene Set Enrichment Analysis (ssGSEA) concepts. Key pathways included:
\begin{itemize}
    \item Proliferation (e.g., \textit{MKI67, targets})
    \item Angiogenesis (e.g., \textit{VEGFA})
    \item Apoptosis
    \item Estrogen Response
    \item Invasion/EMT
\end{itemize}
For the validation cohort, clinical features (Age, Tumor Size, Lymph Node Status) were integrated to enhance predictive power.

\subsection{Model Training}
We trained three machine learning classifiers: Elastic Net (Logistic Regression with L1/L2 regularization), Random Forest, and Gradient Boosting (XGBoost). All models were validated using 5-fold stratified cross-validation.

\section{Results}

\subsection{Predictive Performance}
The Gradient Boosting model achieved superior performance on the external validation set, demonstrating robust generalization.

\begin{table}[h]
\centering
\begin{tabular}{lccc}
\toprule
\textbf{Dataset} & \textbf{Best Model} & \textbf{AUC} & \textbf{Accuracy} \\
\midrule
TCGA (Training) & Elastic Net & 0.81 & 79\% \\
GSE96058 (Validation) & Random Forest & 0.86 & 86\% \\
\bottomrule
\end{tabular}
\caption{Model performance summary.}
\end{table}

\subsection{Explainability}
SHAP (SHapley Additive exPlanations) analysis revealed that the most influential predictors were \textbf{Ki67 status} (Proliferation) and \textbf{Age}. Importantly, high activity of the \textit{Estrogen Response} pathway was associated with lower risk, while high \textit{Proliferation} and \textit{Angiogenesis} scores increased risk, aligning with established cancer biology.

\section{Conclusion}
We successfully developed a robust, interpretable machine learning pipeline for breast cancer risk prediction. By transforming raw genomic data into pathway scores, we improved both generalization (AUC 0.86) and biological transparency. This approach provides a framework for integrating AI into precision oncology.

\end{document}
